\begin{lemma}[Local extension of a convex function to a globally convex one]
Let $E$ be a finite-dimensional real inner-product space, let $\Omega\subseteq E$
be open, let $\psi:E\to\R$ be convex on $\Omega$ (so in particular $\Omega$ is
convex), let $x_0\in E$ and let $r>0$ satisfy
$\overline B(x_0,2r)\subseteq\Omega$. Then there is $\psi':E\to\R$ which is
convex on all of $E$ and satisfies
\[
  \psi'(x)=\psi(x)\qquad\bigl(x\in B(x_0,r)\bigr).
\]
\end{lemma}
\begin{proof}
\emph{Step 1: a bound for $\psi$ and for its subgradients on the inner ball.}
The set $\overline B(x_0,2r)$ is closed and bounded in the finite-dimensional
space $E$, hence compact, and it is contained in $\Omega$, on which $\psi$ is
continuous (preamble). So
\[
  M:=\max_{\overline B(x_0,2r)}|\psi|
\]
is a finite nonnegative real number (preamble).

For each $w\in B(x_0,r)$ we have $w\in\Omega$, so by
\noderef{ConvexSubgradientExists} there is $q_w\in E$ with
\begin{equation}\label{eq:cle-sub}
  \psi(w)+\langle q_w,x-w\rangle\le\psi(x)\qquad(x\in\Omega).
\end{equation}
Fix one such $q_w$ for each $w\in B(x_0,r)$. These vectors are bounded: let
$e\in E$ with $|e|=1$ and put $x=w+re$. Then
$|x-x_0|\le|w-x_0|+r<r+r=2r$, so $x\in\overline B(x_0,2r)\subseteq\Omega$, and
\eqref{eq:cle-sub} gives
\[
  r\langle q_w,e\rangle\le\psi(w+re)-\psi(w)\le M-(-M)=2M .
\]
Dividing by $r>0$ gives $\langle q_w,e\rangle\le2M/r$ for every unit vector $e$,
hence
\begin{equation}\label{eq:cle-qbound}
  |q_w|\le\frac{2M}{r}\qquad\bigl(w\in B(x_0,r)\bigr)
\end{equation}
(take $e=q_w/|q_w|$ when $q_w\ne0$; the bound is trivial otherwise).

\emph{Step 2: the extension.}
For $w\in B(x_0,r)$ let $a_w:E\to\R$ be the affine function
$a_w(x)=\psi(w)+\langle q_w,x-w\rangle$, and define
\[
  \psi'(x)=\sup\bigl\{a_w(x):w\in B(x_0,r)\bigr\}\qquad(x\in E).
\]
The index set is nonempty, since $r>0$ gives $x_0\in B(x_0,r)$. The supremum is
finite for every $x$: by $|\psi(w)|\le M$ and \eqref{eq:cle-qbound},
\[
  a_w(x)\le M+\frac{2M}{r}|x-w|\le M+\frac{2M}{r}\bigl(|x-x_0|+r\bigr)
  \qquad\bigl(w\in B(x_0,r)\bigr),
\]
a bound independent of $w$. Hence $\psi'(x)$ is a real number, it is an upper
bound for every $a_w(x)$, and it is the least such.

\emph{Step 3: $\psi'$ is convex on $E$.}
The set $E$ is convex. Let $x_1,x_2\in E$ and $t_1,t_2\ge0$ with $t_1+t_2=1$, and
let $w\in B(x_0,r)$. Since $a_w$ is affine,
\[
  a_w(t_1x_1+t_2x_2)=t_1a_w(x_1)+t_2a_w(x_2)\le t_1\psi'(x_1)+t_2\psi'(x_2),
\]
using $a_w(x_i)\le\psi'(x_i)$ and $t_i\ge0$. The right-hand side does not depend
on $w$, so it is an upper bound for $\{a_w(t_1x_1+t_2x_2):w\in B(x_0,r)\}$, whose
least upper bound is $\psi'(t_1x_1+t_2x_2)$. Therefore
$\psi'(t_1x_1+t_2x_2)\le t_1\psi'(x_1)+t_2\psi'(x_2)$, i.e. $\psi'$ is convex on
$E$.

\emph{Step 4: $\psi'=\psi$ on $B(x_0,r)$.}
Let $x\in B(x_0,r)$; then $x\in\Omega$. For every $w\in B(x_0,r)$,
\eqref{eq:cle-sub} applied at this $x$ gives $a_w(x)\le\psi(x)$, so $\psi(x)$ is
an upper bound for the set defining $\psi'(x)$ and therefore
$\psi'(x)\le\psi(x)$. Conversely the choice $w=x$ is admissible and gives
$a_x(x)=\psi(x)$, an element of that set, so $\psi'(x)\ge\psi(x)$. Hence
$\psi'(x)=\psi(x)$.
\end{proof}
